\subsection*{執筆時の注意}

\begin{itemize}
  \item \texttt{width} は基本的に論理qubit数またはbinary variable規模として扱う。
  \item \texttt{depth} は定義が論文ごとに異なるため、\texttt{depth definition} を必ず併記する。
  \item ARLは主指標ではなく、width/depthが得られた技術段階を示す補助分類とする。
  \item EV stockと単一VRPインスタンスの車両数Kを直接同一視しない。
  \item 量子優位を示した、と言わない。
  \item 小規模QAOA実験の成功を、実物流規模の有効性へ直接外挿しない。
\end{itemize}

\subsection*{次に本文へ入れる材料}

\begin{itemize}
  \item \texttt{social\_stage\_sources.csv} から社会段階定義の根拠を追加する。
  \item \texttt{technical\_stage\_sources.csv} から技術段階定義の根拠を追加する。
  \item \texttt{roadmap\_gap\_table.md} の各gapを本文用の短い段落へ変換する。
  \item Figure 1、Table Dのキャプションを投稿先フォーマットに合わせて整える。
\end{itemize}
